\section{Switching a past and a future with controlled volume}
\label{sec:switching}

The finite-interval estimates now give a geometric operation on pairs of
nearby points.  It joins the future of one point to the past of the other.
We must control the distribution of the resulting point as well as its
orbit: this requires constructing both complementary joins and estimating
their joint Jacobian.

The joining geometry is related to the canonical local-product
coordinates of hyperbolic dynamics \cite[\S I.7]{Smale1967}.
Solving a linear difference equation and then a nonlinear contraction
also has a standard shadowing precedent
\cite[Lemma~3.2 and Proposition~4.1]{ChowLinPalmer1989}.
Here the slow group may include zero or positive rates; we prove the
finite-interval boundary inverse directly. The measure inequality
requires the further determinant and injectivity arguments below.

\begin{theorem}[Smooth switching]\label{thm:switching}
Let $M$ be a compact smooth manifold of positive finite dimension, let
$\nu$ be a smooth positive probability volume in the density sense of
\Cref{sec:introduction}, and let $S:M\to M$
be a smooth diffeomorphism preserving $\nu$.  Fix a Riemannian distance
$d_M$ and $\eta>0$.  There exist positive numbers $B_k$ with
$\sum_k B_k^{16}<\infty$ and a constant $C$ such that the following
holds for all sufficiently large $k$, with $n=2^k$.

There is a finite measurable partition $\mathcal D_k$ of $M$ with
$\#\mathcal D_k\le \exp(CB_kn)$.  Define the probability measure
\begin{equation}\label{eq:pair-law}
 \rho_k=\sum_{\substack{Q\in\mathcal D_k\\\nu(Q)>0}}
       \frac{(\nu|_Q)\otimes(\nu|_Q)}{\nu(Q)}.
\end{equation}
Thus one first samples a cell with its $\nu$-weight and then samples
two independent points in that cell.  There are a measurable set
$G_k\subset M\times M$, with $\rho_k(G_k)\ge1-\eta$, and a measurable
map $z_k:G_k\to M$ such that
\begin{equation}\label{eq:switch-law}
 (z_k)_*(\rho_k|_{G_k})\le 2\nu
\end{equation}
and
\begin{equation}\label{eq:switch-shadow}
 \sup_{(x,y)\in G_k}\left[
 \max_{0\le t\le n}d_M(S^t z_k(x,y),S^t x)
 +\max_{-n\le t\le0}d_M(S^t z_k(x,y),S^t y)
 \right]\longrightarrow0.
\end{equation}
The measure on the left of \eqref{eq:switch-law} is not normalized
after restriction to $G_k$.  The same conclusions hold if $M$ has a
smooth boundary.
\end{theorem}

The proof occupies the rest of this section.  Its first two steps produce
two complementary switches $(z,w)$.  The last two steps show that their
pair map is injective and has volume Jacobian $1+o(1)$.  This is what
will give \eqref{eq:switch-law}.

\subsection{Grids, retained pairs, and fixed endpoint projections}

We may assume $0<\eta<1$.  Apply \Cref{prop:good-frames} with parameter
$\eta/2$.  It supplies $b=b_k\ge n^{-1/10}$,
$\sum_k b_k^2<\infty$, and a compact set of good base points,
denoted here by $\mathcal A_k$ to distinguish it from the retained pair set,
with suitable frame paths on $[-n,n]$, such that
$\nu(M\setminus\mathcal A_k)<\eta/100$.  Set
\begin{equation}\label{eq:switch-scales}
 B=b^{1/8},\quad h_0=e^{-Bn},\quad r=e^{-2Bn},\quad
 L=\lceil b^{1/4}n\rceil,\quad g=b^{1/2},\quad
 \delta=e^{-3Bn/2}.
\end{equation}
Here and below constants depend only on the fixed smooth data, dimension,
and $\eta$.  They never depend on a point, bin, frame witness, or $k$.
In particular,
\begin{equation}\label{eq:switch-scale-order}
 \log n=o(bn),\quad bn=o(gL),\quad gL=o(gn),\quad
 gn=o(L),\quad L=o(Bn).
\end{equation}
These relations absorb every fixed power of a quantity bounded by
$e^{Cbn}$ or $e^{Cgn}$ at the places where it is used below.

Choose a finite system of coordinate patches with smaller coordinate
boxes covering $M$, each smaller box having closure in its patch.
Assign points to the first smaller box containing them.  This gives
disjoint chart regions whose boundaries lie in finitely many smooth
coordinate faces.  The coordinate maps, their inverses, and the chart
density of $\nu$ have uniform bounds on slightly larger compact patches;
the density has a uniform positive lower bound.  Within each region
use the grid of coordinate cubes of side $h_0$, keeping only cubes whose
closures lie in that region.  We call these cubes \emph{coarse bins}.
The discarded part has volume $O(h_0)$: it is contained in a fixed
multiple of the $h_0$-neighborhood of the finitely many region boundaries.
All comparisons at a common starting or endpoint bin use its one fixed
coordinate map.

For a starting bin $Q$ meeting $\mathcal A_k$, \Cref{lem:bin-spaces}
gives a common threshold $\gamma_Q\in[g,Cg]$, a fast dimension $p_Q$,
and reference spaces $E_Q^s,E_Q^f$ at time zero.  The spaces of every
good frame path beginning in $Q$ are within $e^{-cgL}$ of these spaces.
Their complementary projections and the inverse of a basis orthonormal
within each space have norm at most $e^{Cbn}$.  The threshold and the
reference spaces are chosen once for this bin.  This involves finitely
many choices at each scale, not a choice of a frame as a measurable
function on $\mathcal A_k$.

In the starting chart let $E_Q:\R^{d-p_Q}\times\R^{p_Q}\to\R^d$
be the matrix with these orthonormal slow and fast columns, respectively.
Tile the chart by the parallelepipeds $E_Q([0,r)^d+rj)$, $j\in\Z^d$,
and keep the ones whose closures lie in $Q$.  Translations of the grid
are immaterial.  These are the \emph{product cells}.  Although
$\|E_Q^{-1}\|\le e^{Cbn}$, the forward norm satisfies
$\|E_Q\|\le\sqrt2$: for $(a,b)$,
$|E_Q(a,b)|\le |a|+|b|$.  Thus every cell has physical diameter at
most $C_dr$ and volume at least $e^{-Cbn}r^d$ up to a fixed factor.
The total number of full cells is consequently at most
\begin{equation}\label{eq:cell-count}
 C e^{Cbn}r^{-d}\le e^{CBn}.
\end{equation}
Every point of a coarse bin at distance greater than $C_dr$ from its
walls belongs to a full product cell.  Hence the partial-cell loss,
summed over all coarse bins, is $O(r/h_0)=o(1)$.
Use one additional cell for all omitted points and for bins not meeting
$\mathcal A_k$.  Assign grid boundaries arbitrarily; they are null.
This completes the partition $\mathcal D_k$ without changing
\eqref{eq:cell-count}.

For two points $x=(x_s,x_f)$ and $y=(y_s,y_f)$ in the product
coordinates of one cell, the ideal exchange
\[
 z_0=(y_s,x_f),\qquad w_0=(x_s,y_f)
\]
stays in that cell.  These points combine the slow coordinate of one
input with the fast coordinate of the other.  For the pairs retained
below, we will construct actual orbit joins $z,w$ whose product
coordinates differ from $z_0,w_0$ by $o(r)$.
Removing thin coordinate strips will therefore keep those
joins in their original cell.  This is why the cells follow the two
reference spaces, which need not be orthogonal; see
\Cref{fig:ideal-cell-exchange}.

\begin{figure}[htbp]
\centering
\begin{tikzpicture}[x=.9cm,y=.9cm,>=Stealth,
 point/.style={circle,fill,inner sep=1.6pt}]
 \draw[thick] (0,0)--(4.8,0)--(6.3,2.8)--(1.5,2.8)--cycle;
 \draw[dashed,gray] (1.104,0)--(2.604,2.8);
 \draw[dashed,gray] (3.504,0)--(5.004,2.8);
 \draw[dashed,gray] (.39,.728)--(5.19,.728);
 \draw[dashed,gray] (1.05,1.96)--(5.85,1.96);
 \node[point,label=above left:{$x$}] at (2.154,1.96) {};
 \node[point,label=below right:{$y$}] at (3.894,.728) {};
 \node[point,label=above right:{$z_0$}] at (4.554,1.96) {};
 \node[point,label=below left:{$w_0$}] at (1.494,.728) {};
 \draw[->] (0,-.4)--(1.5,-.4) node[midway,below] {slow};
 \draw[->] (-.45,0)--(.05,.933) node[midway,above left] {fast};
\end{tikzpicture}
\caption{A two-dimensional schematic of the ideal exchange in a
product cell.  Dashed lines hold one product coordinate fixed.
The two coordinate directions may be oblique.  For retained pairs, the
actual joins are proved to differ from $z_0,w_0$ by $o(r)$ in
\Cref{lem:switch-orbits}.}
\label{fig:ideal-cell-exchange}
\end{figure}

Fix a small number $\tau>0$, independent of $k$.  The interior of a
product cell retained at this stage consists of points whose every
product coordinate is at least $\tau r$ from its two cell faces.
The discarded coordinate strips have at most $2d\tau$ of its Euclidean
volume.  The constant determinant of $E_Q$ cancels in this fraction.
Moreover, the ratio between the largest and smallest smooth density on
the physical cell is $1+O(r)$, uniformly in $Q$.  Thus the total smooth
volume of these strips is at most $2d\tau+o(1)$.  The corresponding
strips of width $\tau h_0$ in coarse bins have the same bound.  These
are estimates for the original measure $\nu$, without conditioning on
membership in $\mathcal A_k$.

If $M$ has boundary, also exclude points having any iterate at an integer
time in $[-n,n]$ at distance at most $c_0h_0$ from $\partial M$, where
$c_0>0$ is fixed.  A smooth boundary collar has volume $O(h_0)$, so
invariance bounds this loss by $O(nh_0)=o(1)$.  The latter limit follows
from $Bn\ge n^{79/80}$.  Local coordinate extensions across the boundary
may be used for uniform derivative bounds; all orbit tubes constructed
below lie in the original manifold.

Let $R_k$ be the Borel set of points in $\mathcal A_k$ that lie in the
retained interiors of full product cells, whose two endpoints $S^{\pm n}u$
lie in coarse bins at coordinate distance at least $\tau h_0$ from
their faces, and that satisfy the boundary provision when needed.
Bins not meeting $\mathcal A_k$ contain no such point.  The preceding
estimates and invariance give
\begin{equation}\label{eq:retained-point-loss}
 \nu(M\setminus R_k)\le \eta/100+C_d\tau+o(1).
\end{equation}
Choose $\tau$ so that $C_d\tau<\eta/100$, and then increase $k$ so
that the right-hand side is less than $\eta/4$.
Both marginals of \eqref{eq:pair-law} are exactly $\nu$.  Consequently,
for
\begin{equation}\label{eq:retained-pairs}
 G_k=(R_k\times R_k)\cap
       \bigcup_{Q'\text{ a full product cell}}(Q'\times Q'),
 \qquad \rho_k(G_k)\ge1-\eta/2\ge1-\eta.
\end{equation}
The measure $\rho_k$ is supported on pairs in the same partition cell,
and $R_k$ meets only full product cells.  Hence
$\rho_k(G_k)=\rho_k(R_k\times R_k)$, and the two marginal bounds give
the displayed estimate.  No uniform distribution of good points inside
individual cells is required.

We next fix the maps that will specify the switch.  For each starting
bin $Q$ and endpoint bin $U$ reached at time $n$ by a point of
$\mathcal A_k\cap Q$, choose a representative endpoint fast space.
Let $P^+_{Q,U}:\R^d\to\R^{p_Q}$ be its orthogonal projection, followed
by a fixed orthonormal identification with $\R^{p_Q}$.  Similarly, for
an endpoint bin $V$ reached at time $-n$, define
$P^-_{Q,V}:\R^d\to\R^{d-p_Q}$ from a representative slow space.
By the endpoint conclusion of \Cref{lem:bin-spaces}, these projections
restrict to isomorphisms on the actual fast and slow endpoint spaces
of every corresponding good path, with inverse norm at most $2$ for
large $k$.  The maps depend only on the displayed bins and the starting
threshold.  They do not depend on a later choice of input in those bins.
Zero-dimensional projections and determinants have their usual empty
space meanings.

For a retained pair $(x,y)$, write $Q$ for its coarse starting bin and
$U_x,V_x,U_y,V_y$ for the coarse bins of
$S^nx,S^{-n}x,S^ny,S^{-n}y$, respectively.  If $\chi_U$ denotes the
coordinate map attached to a bin, define the following smooth maps on
the open sets where the coordinates make sense:
\begin{equation}\label{eq:endpoint-maps}
 \begin{aligned}
 F_x(u)&=P^+_{Q,U_x}\chi_{U_x}(S^nu),&
 P_x(u)&=P^-_{Q,V_x}\chi_{V_x}(S^{-n}u),\\
 F_y(u)&=P^+_{Q,U_y}\chi_{U_y}(S^nu),&
 P_y(u)&=P^-_{Q,V_y}\chi_{V_y}(S^{-n}u).
 \end{aligned}
\end{equation}
The subscripts denote fixed bin labels.  We seek $z,w$ satisfying
\begin{equation}\label{eq:four-switch-equations}
 F_x(z)=F_x(x),\quad P_y(z)=P_y(y),\qquad
 F_y(w)=F_y(y),\quad P_x(w)=P_x(x).
\end{equation}
The next step both solves these equations near the prescribed halves
and gives a larger neighborhood in which each solution is unique.
\Cref{fig:switch-data} displays the required exchange of endpoint data.

\begin{figure}[htbp]
\centering
\begin{tikzpicture}[>=Stealth,
 point/.style={circle,draw,minimum size=7mm,inner sep=1pt},
 future/.style={->,thick,blue!65!black},
 past/.style={->,thick,dashed,orange!75!black}]
 \node[point] (x) at (-3,0.85) {$x$};
 \node[point] (y) at (-3,-0.85) {$y$};
 \node[point] (z) at (3,0.85) {$z$};
 \node[point] (w) at (3,-0.85) {$w$};
 \draw[future] (x) -- node[above] {$F_x$} (z);
 \draw[future] (y) -- node[below] {$F_y$} (w);
 \draw[past] (x) -- node[pos=0.27,above,sloped] {$P_x$} (w);
 \draw[past] (y) -- node[pos=0.27,below,sloped] {$P_y$} (z);
\end{tikzpicture}
\caption{The endpoint equations for the complementary switches.
Solid arrows copy projected future coordinates; dashed arrows copy
projected past coordinates.  The arrows represent endpoint data,
not trajectories in the manifold.}
\label{fig:switch-data}
\end{figure}

\subsection{A boundary inverse and a nonlinear orbit}

\begin{lemma}\label{lem:switch-orbits}
For every sufficiently large $k$ and every $(x,y)\in G_k$, the first
two equations of \eqref{eq:four-switch-equations} have a solution $z$
with
\begin{equation}\label{eq:sharp-shadow}
 \max_{0\le t\le n}d_M(S^tz,S^tx)
 +\max_{-n\le t\le0}d_M(S^tz,S^ty)
 \le e^{Cgn}r.
\end{equation}
In fixed regular coordinates along these halves the solution is unique
among displacement arrays of sup norm at most $\delta$.
The complementary solution $w$ has the analogous properties with $x,y$
exchanged.  In the product coordinates of the common cell, if
$x=(x_s,x_f)$ and $y=(y_s,y_f)$, then
\begin{equation}\label{eq:ideal-switch}
 z=(y_s,x_f)+o(r),\qquad w=(x_s,y_f)+o(r),
\end{equation}
uniformly over all retained pairs.  Both outputs remain in the original
product cell, and their relevant endpoints remain in the respective
four original endpoint bins.
\end{lemma}

\begin{proof}
\emph{Uniform coordinates.}
At each reference location choose one of finitely many slightly larger
regular patches containing a fixed-radius neighborhood of that location.
At time zero both strings use $\chi_Q$; at the terminal times they use
the endpoint charts already specified.  Shrinking the fixed radius once,
the one-step coordinate maps
$\chi_{t+1}\circ S\circ\chi_t^{-1}$ and their first two derivatives
are uniformly bounded wherever they will be used.  The inverse
one-step maps also have uniform first and second derivative bounds.
For example, differentiating the inverse identity bounds their second
derivatives by a fixed multiple of the forward second-derivative bound
and the cube of the inverse first-derivative bound.  This
follows from a finite atlas with nested compact patches and smoothness
on compact $M$; it does not depend on the number of steps.  A patch
can be chosen separately at each internal time.  No iterate $S^n$
is assigned a uniform second derivative bound.

Use displacement variables $v_t^+$ about $S^tx$, $0\le t\le n$,
and $v_t^-$ about $S^ty$, $-n\le t\le0$.  Write $A_t^+,A_t^-$
for the derivative steps in these charts.  Matching the two locations
at zero and imposing the terminal equations gives the linear system
\begin{equation}\label{eq:linear-boundary-system}
 \begin{aligned}
 v_{t+1}^{\pm}-A_t^{\pm}v_t^{\pm}&=f_t^{\pm},\\
 v_0^+-v_0^-&=m,\\
 P^+_{Q,U_x}v_n^+&=h_+,\qquad
 P^-_{Q,V_y}v_{-n}^-=h_-.
 \end{aligned}
\end{equation}
The matching datum for the actual problem is
$m=\chi_Q(y)-\chi_Q(x)$; its norm is $O(r)$.
The actual terminal data are zero.

\emph{The inverse in weighted coordinates.}
Let $\gamma=\gamma_Q$ and put $\widetilde v_t=e^{-\gamma t}v_t$.
Use the same weighting on equation data, with the step datum at its
output time $t+1$.  By \Cref{prop:finite-splitting}, the weighted
propagators on slow spaces forward and fast spaces backward have norm
at most $e^{Cbn}e^{-g|t-s|}$.  The associated projections have norm
at most $e^{Cbn}$.  Geometric sums cost at most
$C/g\le e^{Cbn}$, by \eqref{eq:switch-scale-order} and the floor on $b$.

On the future half, prescribe a slow vector $a\in E_0^s(x)$.
Variation of constants gives its slow part at every time by propagating
$a$ forward and adding the projected step data.  If $s_n$ is the
resulting slow value at time $n$, the fast value there is uniquely
\[
 c_n=(P^+_{Q,U_x}|_{E_n^f(x)})^{-1}
                (\widetilde h_+-P^+_{Q,U_x}s_n).
\]
Propagate $c_n$ and the fast step data backward.  All forced terms have
norm at most $e^{Cbn}$ times the sup norm of the weighted data.
The dependence of the resulting fast value at zero on $a$ is the map
\begin{equation}\label{eq:future-kernel-graph}
 K_+=-\widetilde A_x(0,n)|_{E_n^f(x)}
       (P^+_{Q,U_x}|_{E_n^f(x)})^{-1}
       P^+_{Q,U_x}\widetilde A_x(n,0)|_{E_0^s(x)},
 \qquad \|K_+\|\le e^{-2gn+Cbn}.
\end{equation}
Here $\widetilde A(t,s)=e^{-\gamma(t-s)}A(t,s)$.
Both decaying factors are present in \eqref{eq:future-kernel-graph};
there is no assumption that the endpoint projection annihilates the
slow space.  Thus $\widetilde v_0^+=a+K_+a+q_+$, with $q_+$ bounded
by the weighted data as above.

On the past half prescribe $b'\in E_0^f(y)$, propagate fast backward,
use the slow endpoint projection to fix the slow value at $-n$, and
then propagate slow forward.  This gives
$\widetilde v_0^-=b'+K_-b'+q_-$, where
\begin{equation}\label{eq:past-kernel-graph}
 \|K_-\|\le e^{-2gn+Cbn},\qquad
 \|q_-\|\le e^{Cbn}\|\text{weighted data}\|_{\infty}.
\end{equation}
The map $(a,b')\mapsto a-b'$ from $E_0^s(x)\oplus E_0^f(y)$ to
the starting coordinate space has inverse norm at most $e^{Cbn}$.
Indeed the corresponding reference spaces are complementary with this
bound, and their perturbation is $e^{-cgL}$; multiplication by the
reference inverse still gives a perturbation tending to zero.
The terms $K_+a-K_-b'$ are smaller still.  A Neumann series therefore
solves the matching equation with inverse cost $e^{Cbn}$.
This proves an $e^{Cbn}$ inverse bound for the full weighted system.
Removing the weights costs at most $e^{2\gamma(n+1)}$, so its
unweighted inverse $\mathcal L^{-1}$ satisfies
\begin{equation}\label{eq:boundary-inverse}
 \|\mathcal L^{-1}\|_{\infty\to\infty}\le Q_k,
 \qquad Q_k=e^{Cgn}.
\end{equation}
If a group is empty, omit its variables, propagation formula, and
terminal datum.  The matching map on the remaining spaces and the
same bounds apply, including the cases $p_Q=0,d$.

\emph{Contraction with a larger uniqueness radius.}
Taylor expansion of each one-step map gives a residual
$\mathcal R(v)$ with
\begin{equation}\label{eq:one-step-remainder}
 \|\mathcal R(v)\|_\infty\le C\|v\|_\infty^2,
 \qquad
 \|\mathcal R(v)-\mathcal R(u)\|_\infty
       \le C\delta\|v-u\|_\infty
\end{equation}
on the sup norm ball of radius $\delta$.  Every residual coordinate
depends on just its one step, so the constant has no factor $n$.
The matching and terminal equations are affine and have zero nonlinear
residual.  The nonlinear problem is
$v=\mathcal L^{-1}(a_0+\mathcal R(v))$, where
$\|a_0\|_\infty\le Cr$.
The scale relations give
\[
 CQ_k\delta\longrightarrow0,\qquad Q_kr/\delta\longrightarrow0.
\]
For large $k$ the displayed map sends the $\delta$-ball into its
$\delta/2$-ball and has Lipschitz constant at most $1/4$.
It has a unique fixed point in the larger ball, with
$\|v\|_\infty\le 2CQ_kr$.  The exact one-step equations and the
matching equation make this array one genuine orbit; call its point
at zero $z$.  Chart bounds give \eqref{eq:sharp-shadow}.
In the boundary case its distance from each reference orbit is
$o(h_0)$, so it stays inside $M$.

\emph{Time-zero coordinates and wall margins.}
For the linear solution with zero step and terminal data, the two
initial displacement spaces are the graphs in
\eqref{eq:future-kernel-graph} and \eqref{eq:past-kernel-graph}.
In the reference product coordinates they differ from the slow and
fast coordinate spaces by at most
$e^{Cbn}(e^{-cgL}+e^{-2gn+Cbn})$.  Solving their matching equation
therefore differs from ideal exchange of the fast and slow coordinates
by at most
\begin{equation}\label{eq:linear-switch-error}
 e^{Cbn}\bigl(e^{-cgL}+e^{-2gn+Cbn}\bigr)r=o(r).
\end{equation}
All basis conversion and matching inverse factors are included in the
constant $C$.  The nonlinear solution differs from the linear one by
at most $CQ_k(2CQ_kr)^2=e^{Cgn}r^2$ in ordinary coordinates, and still
by $e^{Cgn}r^2=o(r)$ in product coordinates.  This proves
\eqref{eq:ideal-switch}.  The proof for $w$ is identical with the
two inputs exchanged.

Increase $k$ so that the error in each product coordinate is less than
$\tau r/4$.  The ideal exchanged point has distance at least $\tau r$
from every product face because both inputs do.  Thus $z,w$ have
coordinate margins at least $3\tau r/4$ in their original cell.
Similarly, \eqref{eq:sharp-shadow} in the fixed endpoint charts is
less than $\tau h_0/4$, so each corresponding endpoint retains margin
at least $3\tau h_0/4$ in its original coarse bin.  These strict
margins will recover all labels from the outputs.
The point obtained does not depend on the auxiliary internal charts or
on the frame witnesses used to justify the estimates.  Two such
constructions have the same endpoint equations and sharp radius
$e^{Cgn}r=o(\delta)$.  In either one's charts both are therefore in
its uniqueness ball, and hence coincide.
\end{proof}

\subsection{Relative minors and the joint determinant}

We have constructed complementary orbit joins.  A single join need not
have a convenient Jacobian, so we now keep both of them.  Their endpoint
coordinates are permuted, which will cancel the large expansion and
contraction factors.

For the fixed labels of a retained pair define
\begin{equation}\label{eq:joint-endpoint-map}
 H_{ab}(u)=(F_a(u),P_b(u)),\qquad a,b\in\{x,y\}.
\end{equation}
Each target has dimension $p_Q+(d-p_Q)=d$ and uses the fixed
orthonormal coordinates in the projection ranges.

\begin{lemma}\label{lem:switch-jacobian}
All four derivatives $DH_{xx}(x)$, $DH_{yy}(y)$,
$DH_{xy}(z)$, and $DH_{yx}(w)$ are nonsingular.  Locally the equations
\eqref{eq:four-switch-equations} define a smooth pair map
$\Phi(x,y)=(z,w)$ with
\begin{equation}\label{eq:pair-jacobian}
 J_{\nu\otimes\nu}\Phi(x,y)=1+o(1),
\end{equation}
uniformly on the retained pairs.  The conclusion includes empty slow
or fast groups.
\end{lemma}

\begin{proof}
Put $p=p_Q$.  For the future map at $x$, use the actual fast and slow
spaces first, with orthonormal bases within each.  The fast $p$-minor
has absolute value at least $e^{\Lambda_x n-Cbn}$, where
$\Lambda_x$ is the sum of its fast rates.  This follows from the
fast volume lower bound in \Cref{prop:finite-splitting} and from the
uniformly invertible restriction of the endpoint projection.
Every other maximal minor contains a slow starting vector and is at
most $e^{(\Lambda_x-2g)n+Cbn}$ when both groups occur.
Replacing the actual bases with the common reference basis changes
their spaces by $e^{-cgL}$.  After a change of orthonormal basis within
each space, multilinearity and the full exterior norm bound show that
the leading minor stays nonzero and that all other maximal minors
are $o(1)$ times its absolute value.  Factors from the inverse of the
possibly oblique starting basis are at most $e^{Cbn}$.  Call the
resulting leading absolute minor $A_x$.  More explicitly, when
$0<p<d$, every nonleading maximal future minor $\Delta$ in the common
reference basis satisfies
\[
 \frac{|\Delta|}{A_x}
 \le C\left(e^{-2gn+Cbn}+e^{-cgL+Cbn}\right)=o(1).
\]
The first term comes from the mixed-exterior estimate.  The second
comes from the change of spaces, using the full exterior upper bound
and the leading-minor lower bound.  Increasing $C$ includes the
oblique-basis factors.  If $p=0$ or $p=d$, there are no nonleading
maximal future minors, so this comparison is unnecessary.

Define $A_y$ in the same way for the future of $y$.
For the past maps the leading
columns are slow, and the backward exterior rate is
\[
 \Omega_x=-\sum_{\text{slow at }x}\lambda_j(x).
\]
Denote their leading absolute minors by $B_x,B_y$.
They have lower bounds $e^{\Omega_x n-Cbn}$ and
$e^{\Omega_y n-Cbn}$, respectively, and all other maximal minors are
relatively negligible.  No sign is prescribed for $\Lambda_x$ or
$\Omega_x$ in a perturbation estimate.

We must transfer the future minors from $x$ to $z$ relative to $A_x$,
even if this leading quantity is small.  Use the reference charts
along the future of $x$ for both nearby orbits.  By
\eqref{eq:sharp-shadow}, their one-step derivatives differ by at most
$C\delta$ for large $k$.  On every unperturbed subinterval the
$p$-th exterior product has norm at most
$e^{\Lambda_x\ell+Cbn}$.
Applying \Cref{lem:relative-exterior} gives, for every fixed $C_0$,
\begin{equation}\label{eq:relative-minor-transfer}
 \left\|\bigwedge^pD(S^n)_z-
                   \bigwedge^pD(S^n)_x\right\|
       =o\bigl(e^{\Lambda_xn-C_0bn}\bigr),
\end{equation}
where the derivatives are represented in these common starting and
endpoint coordinates.  To see the scale in this application, an
expansion term with $q\ge1$ altered factors is bounded by
\[
 e^{\Lambda_x(n-q)+(q+1)C_1bn}(C\delta)^q.
\]
The binomial sum, divided by $e^{\Lambda_xn-C_0bn}$, tends to zero
because $n\delta e^{Cbn}\to0$.  The bounded factor
$e^{q|\Lambda_x|}$ covers either sign of the rate.  The same reasoning
backward uses $\Omega_y$ and $e^{q|\Omega_y|}$.
Fixed endpoint projections and the common starting basis preserve
\eqref{eq:relative-minor-transfer}, after increasing $C_0$ if an
inverse basis factor is needed.  Hence at $z$ the leading future and
past minors are $A_x(1+o(1))$ and $B_y(1+o(1))$, with the same relative
smallness of all the other minors.  At $w$ they are $A_y$ and $B_x$.

Expand each full determinant along its $p$ future rows.  There is one
leading product, from fast future columns and slow past columns.
Every other term has a relatively negligible maximal minor.  In the
common starting basis the four absolute determinants are therefore
\begin{equation}\label{eq:four-leading-products}
 \begin{array}{c|cccc}
 \text{map and point}&H_{xx}(x)&H_{yy}(y)&H_{xy}(z)&H_{yx}(w)\\ \hline
 \text{absolute determinant}&
 A_xB_x(1+o(1))&A_yB_y(1+o(1))&
 A_xB_y(1+o(1))&A_yB_x(1+o(1)).
 \end{array}
\end{equation}
All leading factors are positive.  In particular these maps are
locally invertible.  If $p=0$, future targets and minors are empty and
$A_x=A_y=1$; if $p=d$, the past minors are $B_x=B_y=1$.  The same
determinants and conclusions apply, without mixed minors of the
missing kind.

The defining equations say that
$(H_{xy}(z),H_{yx}(w))$ is a fixed permutation of
$(H_{xx}(x),H_{yy}(y))$.  Differentiation gives the exact identity
\begin{equation}\label{eq:exact-switch-determinant}
 |\det D\Phi(x,y)|=
 \frac{|\det DH_{xx}(x)|\,|\det DH_{yy}(y)|}
 {|\det DH_{xy}(z)|\,|\det DH_{yx}(w)|}.
\end{equation}
Each full derivative represented in the common product basis differs
from its coordinate derivative by the factor $\det E_Q$.  There are
two such factors in the numerator and two in the denominator, so
they cancel exactly.  Substituting \eqref{eq:four-leading-products}
cancels $A_x,B_x,A_y,B_y$ individually, with no comparison between the
rates of $x$ and $y$.  Thus the coordinate Jacobian is $1+o(1)$.
All four starting points lie in the same cell of diameter $O(r)$.
The ratio of the two output smooth densities to the two input densities
is therefore $1+o(1)$ as well.  This proves
\eqref{eq:pair-jacobian}.
\end{proof}

\subsection{Measurability, recovered labels, and global injectivity}

The Jacobian calculation is local.  To convert it into the desired
measure inequality we must ensure that no multiplicity arises from
different choices of bins or local branches.

\begin{lemma}\label{lem:switch-injective}
The map $\Phi:G_k\to M\times M$, $\Phi(x,y)=(z,w)$, is measurable
and injective.  Its restriction to countably many measurable pieces
agrees with smooth local diffeomorphisms having the Jacobian in
\eqref{eq:pair-jacobian}.
\end{lemma}

\begin{proof}
First fix a retained pair.  Its cell and four endpoint labels are
locally constant among retained inputs near it, because of the positive
wall margins.  By \Cref{lem:switch-jacobian}, the implicit equations
define a smooth local pair branch near its outputs.  This branch
agrees with the constructed switch on nearby retained inputs.
Indeed, for this fixed $k$, continuity of the finitely many maps
$S^t$, $|t|\le n$, makes their reference halves arbitrarily close
when their inputs are close.  The constructed solutions have sharp
radius $e^{Cgn}r=o(\delta)$ about those halves.  Both the old local
implicit branch and the construction for the nearby pair therefore
lie in the old $\delta$-tube, and satisfy the same fixed endpoint
equations with the nearby right-hand sides.  The contraction uniqueness
argument identifies them even for these changed data: two solutions
$u,v$ in the old ball satisfy
$\mathcal L(v-u)=\mathcal R(v)-\mathcal R(u)$, so
$\|v-u\|_\infty\le CQ_k\delta\|v-u\|_\infty$ forces equality.
Both candidate arrays here use the old reference halves. Their
matching datum is therefore the same old coordinate difference;
only the prescribed endpoint values change with the nearby inputs.
The charts along the old pair may be kept fixed in this neighborhood.  In particular no regularity of a choice of good
frame witnesses is needed.

The retained set is Borel, and $M\times M$ is second countable.
A countable subcollection of these relative neighborhoods covers it.
Subtracting earlier members partitions it into countably many Borel
pieces.  On each piece $\Phi$ is the restriction of the indicated
smooth local diffeomorphism; this proves measurability and supplies
the pieces for change of variables.

For injectivity suppose
$\Phi(x,y)=\Phi(x',y')=(z,w)$.  The interiors of product cells are
disjoint, and both outputs stayed in the original cell.  Thus the
two input pairs have the same product cell and hence the same starting
bin $Q$, threshold, and dimensions.  The four endpoint-bin labels
are recovered from the ordered output pair by
\begin{equation}\label{eq:recovered-bin-labels}
 (U_x,V_x,U_y,V_y)=
 \bigl(\operatorname{bin}(S^nz),\operatorname{bin}(S^{-n}w),
       \operatorname{bin}(S^nw),\operatorname{bin}(S^{-n}z)\bigr).
\end{equation}
All four endpoints have positive margins in their bins.  Therefore
these labels, and exactly the four projections and coordinate maps in
\eqref{eq:endpoint-maps}, are the same for the competing inputs.

On the future half, both $x$ and $x'$ are within $e^{Cgn}r$ of $z$;
on the past half both are within that distance of $w$.  Consequently
their full orbits on $[-n,n]$ are within $2e^{Cgn}r=o(\delta)$ of
each other.  Moreover the recovered equations give
\[
 F_x(x')=F_x(z)=F_x(x),\qquad
 P_x(x')=P_x(w)=P_x(x).
\]
Apply the boundary uniqueness argument of \Cref{lem:switch-orbits}
with both reference halves equal to the true good orbit of $x$.
It uses $x$'s endpoint projections, zero matching datum, and zero
terminal data.  The orbit of $x$ is one solution.  In the same charts
the orbit of $x'$ is another solution in its $\delta$-ball, since its
distance is $o(\delta)$ even after fixed chart conversions.  Uniqueness
gives $x'=x$.  The same argument gives $y'=y$.
Only the true input orbits were required to be good; neither $z$ nor
$w$ was required to belong to $\mathcal A_k$ or to $R_k$.
\end{proof}

\begin{proof}[Proof of \Cref{thm:switching}]
The construction above applies for all sufficiently large $k$.
Equation \eqref{eq:switch-scales} gives
$\sum_kB_k^{16}=\sum_kb_k^2<\infty$; any finitely many unused terms
can be chosen positively.  The partition bound follows from
\eqref{eq:cell-count}, retention from \eqref{eq:retained-pairs}, and
uniform shadowing from \eqref{eq:sharp-shadow}.

It remains to prove the output measure bound.  On the square of a
positive-volume cell $Q'$, the original law $\rho_k$ has density
$1/\nu(Q')$ with respect to $\nu\otimes\nu$.  The map $\Phi$ sends
retained pairs in that square into the same square.  By
\Cref{lem:switch-jacobian}, its inverse volume Jacobian is at most $2$
for all sufficiently large $k$.  Apply change of variables on the
countably many pieces of \Cref{lem:switch-injective}.  Their images
are disjoint by global injectivity.  Summing therefore gives
\[
 \Phi_*(\rho_k|_{G_k})\le 2\rho_k.
\]
The factor $1/\nu(Q')$ is unchanged.  Restriction to endpoint-label
pieces or to $G_k$ is never followed by normalization, and zero-volume
cells contribute nothing.  Taking first marginals proves
\eqref{eq:switch-law}, because the first marginal of $\rho_k$ is
$\nu$.  The construction already handled the optional smooth boundary.
\end{proof}
